\scanpage{100}
\mspage{\textcircled{26}}
Formulaire faisant intervenir $\overset{\mathbb L}{\otimes}$ et $D^-$

\textcircled{1} Généralités valables sur tout topos annelé

Module \uline{plat} $F$ : $F\otimes\cdot$ est exact.

$F$ plat $\Longleftrightarrow$ Pour tout injectif $I$, $\mathscr Hom(F,I)$ est injectif.

\uline{NB} Dans le cas actuel, il suffit de vérifier sur les fibres.

Tout \deleted{\unsure{f}}Module $F$ \deleted{admet des \unsure{ré}} est quotient d'un plat, donc admet des résolutions plates : on prend les $A_U$ ($U\in\mathrm{Ob}\,X_{\mathrm{ét}}$). Il s'ensuit \inserted{\unsure{formellt|formellmt}} que pour tt $\deleted{\unsure{K}}L_{\bullet}\in\mathrm{Ob}\,D^-(X)$, existe \unsure{liaison|homom.} $L'_{\bullet}\mathbin{-}L_{\bullet}$, avec $L'_{\bullet}$ plat en \unsure{tout degré|tous degré|tous degrés}, qui est un quasi-isomorphisme.

De ceci résulte formellement que $D^-(X)$ est équivalent : la catégorie obtenue en prenant \inserted{la catégorie $K^-(X)^{\mathrm{plat}}$} des complexes $L_{\bullet}$ bornés sup. plats \deleted{\unsure{en}} tous degrés, \deleted{\unsure{et}} en prenant les hom. à homotopie près, puis prenant un quotient par la \deleted{relation} \inserted{famille} des quasi-isomorphismes.

Ceci va nous permettre de prendre \deleted{des} ``foncteur dérivé gauche'' du foncteur $\otimes$, en prenant dans $C^-(X)^{\mathrm{plat}}$ le $\otimes$ ordinaire ; on \unsure{note|voit} que ça passe au quotient par quasi-isomorphismes. Donc si $L_{\bullet},\deleted{L}M_{\bullet}$ sont deux complexes ordinaires $\in\mathrm{Ob}\,D^-(X)$, on \deleted{définit} calcule $L_{\bullet}\overset{L}{\otimes}\deleted{L}M_{\bullet}$ en prenant des résolutions plates $L'_{\bullet}\longrightarrow{}L_{\bullet}$, $M'_{\bullet}\mathbin{-}M_{\bullet}$ et prenant $L'_{\bullet}\otimes M'_{\bullet}$.
% archive sheet number: 99
